%% =====================================================================
%%  Mode d'emploi : la structure d'une unité et les repères du livre
%% =====================================================================
\liminaire{mode}{Mode d'emploi}{Une unité en sept étapes, toujours dans le même ordre}

Les vingt unités du livre sont construites sur le même modèle. Après une page d'ouverture qui
présente la situation, les objectifs et votre programme de la semaine, chaque unité enchaîne
sept rubriques, repérables à leur couleur et à leur pictogramme.\par\vspace{5mm}

\newcommand{\otEtape}[5]{% clé de rubrique, titre, texte, durée, exemple
  \noindent\begin{tikzpicture}[baseline=(t.north)]
    \fill[\otRubCoul{#1},rounded corners=2.4mm] (0,-11mm) rectangle (11mm,0);
    \otDessinIcone{\otRubIco{#1}}{white}{\otRubCoul{#1}}{6.6mm}{(5.5mm,-5.5mm)}
    \node[anchor=north west,inner sep=0pt] (t) at (0,0) {};
  \end{tikzpicture}\hspace{4mm}%
  \begin{minipage}[t]{\dimexpr\linewidth-15mm\relax}
    {\otSemi\fontsize{7.2}{9}\selectfont\color{\otRubCoul{#1}}\otCaps[12]{\otRubEN{#1}}}\hspace{2mm}%
    {\otSemi\fontsize{12}{14}\selectfont #2}\hfill{\otRegular\fontsize{8.4}{10}\selectfont\color{otMuted}#4}\par\vspace{1mm}
    {\fontsize{10.2}{13}\selectfont #3\par}
    \if\relax\detokenize{#5}\relax\else\vspace{1mm}{\fontsize{9.4}{12}\selectfont\color{otMuted}#5\par}\fi
  \end{minipage}\par\vspace{5mm}}

\otEtape{ecouter}{Dialogues}{Deux scènes de la vie de Léa, en anglais avec leur traduction. Chaque
  dialogue existe en \textbf{version lente} et en \textbf{version normale}. Les expressions à retenir
  sont surlignées ; les numéros orange renvoient aux explications « Bien comprendre ».}{30 min}{}
\otEtape{mots}{Vocabulaire}{Les mots et expressions de l'unité, classés par situation, avec leur
  prononciation et leur traduction. Une piste audio pour les écouter et les répéter.}{30 min}{}
\otEtape{grammaire}{Grammaire}{Une ou deux règles, expliquées en français avec des tableaux clairs, des
  exemples et les pièges typiques des francophones.}{30 min}{}
\otEtape{sons}{Prononciation}{Un son difficile pour les francophones, avec la technique pour le
  produire et des paires de mots pour entraîner votre oreille.}{15 min}{}
\otEtape{culture}{Culture}{Les usages du monde anglophone : ce qui se dit, ce qui se fait… et les
  erreurs à éviter.}{15 min}{}
\otEtape{pratique}{Exercices}{Une douzaine d'exercices progressifs pour lire, écrire, écouter et parler.
  Tous sont corrigés et expliqués en fin de livre.}{30 min}{}
\otEtape{bilan}{Bilan}{L'essentiel de l'unité sur une page, une auto-évaluation et un test express
  de dix questions pour vérifier vos acquis.}{30 min}{}

\clearpage
\otIntertitre{Les repères du livre}
\begin{paire}
\begin{regle}[Grammaire]
Les règles à connaître, présentées simplement, avec des exemples.
\end{regle}
\begin{attention}[Pièges]
Les erreurs typiques des francophones :\par\vspace{1mm}
\fauxjuste{I call me Léa.}{My name is Léa.}
\end{attention}
\end{paire}
\begin{paire}
\begin{astuce}[Méthode]
Des conseils concrets pour apprendre plus vite et retenir plus longtemps.
\end{astuce}
\begin{culture}[Usages]
Les codes sociaux et les différences entre le Royaume-Uni et les États-Unis.
\end{culture}
\end{paire}

\begin{tcolorbox}[otBase,colback=white,colframe=otRule,boxrule=0.5pt]
\begin{tabularx}{\linewidth}{@{}>{\centering\arraybackslash}m{30mm}@{\hspace{5mm}}Y@{}}
\otPastillePiste{1.05} & \textbf{Une piste audio.} Le premier nombre est celui de l'unité, le second celui de la
  piste ; puis la durée. Scannez le QR code voisin pour l'écouter.\\[3mm]
\otCompetence{1}{casque}\otCompetence{1}{micro}\otCompetence{1}{livre}\otCompetence{1}{crayon} &
  \textbf{Les compétences travaillées} par un exercice : écouter, parler, lire, écrire.\\[3mm]
\otDifficulte{1}\hspace{3mm}\otDifficulte{3} & \textbf{La difficulté} d'un exercice : de un point (facile) à trois points
  (plus exigeant).\\[3mm]
{\otLight\fontsize{7.5}{9}\selectfont\color{otMuted}corrigé p.~123} & \textbf{Le renvoi au corrigé}, en fin de livre. Les
  corrigés renvoient de même à la page de l'exercice.\\[3mm]
\hl{\enb{Nice to meet you}}\,\note{2} & \textbf{Dans les dialogues}, les expressions à retenir sont surlignées ;
  les numéros renvoient aux explications « Bien comprendre ».\\[3mm]
\ipa{həˈloʊ} & \textbf{La prononciation} en alphabet phonétique, entre barres obliques. Le guide des
  sons, page~\pageref{lim-sons}, vous apprend à la lire.\\
\end{tabularx}
\end{tcolorbox}

\otIntertitre{En fin de livre}
\begin{tcbraster}[raster columns=3,raster equal height=rows,raster column skip=3.5mm,raster before skip=1mm]
\otCartePrg{coche}{otGreenDark}{Corrigés}{Toutes les réponses, avec des explications pas à pas, des « À retenir » et les erreurs fréquentes.}
\otCartePrg{casque}{otNavy}{Transcriptions}{Le texte de toutes les pistes d'exercices, pour vérifier ce que vous avez entendu.}
\otCartePrg{livre}{otOrangeDark}{Lexique}{Tous les mots du tome, de l'anglais vers le français, avec leur prononciation.}
\end{tcbraster}
